\documentclass[11pt,a4paper]{article}

\usepackage[utf8]{inputenc}
\usepackage[T1]{fontenc}
\usepackage{lmodern}
\usepackage[margin=1in]{geometry}
\usepackage{amsmath,amssymb,amsthm,mathtools}
\usepackage{graphicx}
\usepackage{booktabs}
\usepackage{hyperref}
\usepackage{xcolor}
\usepackage{enumitem}
\usepackage{caption}

\hypersetup{
  colorlinks=true,
  linkcolor=blue!50!black,
  citecolor=blue!50!black,
  urlcolor=blue!50!black
}

\title{\bfseries Deutsch-CTC is a limit-cycle theory: \\
numerical evidence for cycle solutions, multi-basin dynamics, \\
and a continuum of chronology-coupling regimes}
\author{Christian Knopp\\\texttt{cknopp@gmail.com}}
\date{2026-05-11}

\begin{document}
\maketitle

\begin{abstract}
We report three numerical findings that revise the standard
fixed-point interpretation of Deutsch's closed-timelike-curve
formulation \cite{Deutsch1991} on small Hilbert spaces. First, the
majority of Clifford-structured D-CTC channels do \emph{not}
converge to a fixed point: they reach a limit cycle of period $k$
predicted exactly by $\arg(\lambda_2(\mathbf{E})) = 2\pi/k$, where
$\mathbf{E}$ is the channel superoperator and $\lambda_2$ its
second-largest eigenvalue. Cycle solutions have $|\lambda_2| = 1$
marginal modes that preserve input information indefinitely, and
their state-distinguisher amplification is $\sim 2\times$ that of
the convergent class. Second, the multi-basin structure of these
cycles means D-CTC dynamics are \emph{not} ergodic across the full
density-matrix space: a single Clifford-class channel admits an
average of 25 distinct limits depending on initial state (vs.\ 1 for
Haar-random channels). Third, the previously-reported binary
``abstract vs.\ direct-CTC'' distinction in chronology-protection
coupling is actually a continuum interpolated by the channel's
direct-CTC content, with the chronology-protection signature
appearing gradually above 50\% direct encoding. Together these
findings imply Deutsch's self-consistency equation should be
re-read as admitting cycle solutions of any period $k \ge 1$, with
the Aaronson--Watrous \cite{AaronsonWatrous2009} state-distinguisher
speedup operating primarily on cycle solutions rather than fixed
points. All code, data, and reproduction scripts are at
\texttt{github.com/Zynerji/systrophe} (v0.17.0).
\end{abstract}

\section{Introduction and motivation}

The Deutsch fixed-point equation for closed-timelike-curve quantum
mechanics \cite{Deutsch1991} reads
\begin{equation}
\rho^* = \mathcal{E}(\rho^*) \quad\text{with}\quad
\mathcal{E}: \rho \mapsto \mathrm{Tr}_{\text{CR}}\!\big[ U(\sigma_{\text{CR}} \otimes \rho) U^\dagger \big]
\label{eq:dctc-fp}
\end{equation}
and is universally treated as defining a fixed-point problem. The
practical algorithm is to iterate $\rho_{n+1} = \mathcal{E}(\rho_n)$
from a random initial state, stopping when $\|\rho_{n+1} - \rho_n\| <
\epsilon$. The Aaronson--Watrous result \cite{AaronsonWatrous2009}
states that an oracle for $\rho^*$ enables polynomial-time PSPACE
computation, including amplifying the distinguishability of two
close mixed states beyond the Helstrom bound \cite{Helstrom1976}.

In a companion paper \cite{Knopp2026DCTC}, we identified the
Clifford-structured channel class (permutation $\times$ diagonal-of-
fourth-roots unitaries) as the empirical site of Aaronson--Watrous
amplification: $\sim 40\%$ of such channels produce purity-$> 0.9$
fixed points, versus $\sim 0\%$ for Haar.

Three deep-dive experiments documented in this note overturn or
sharpen the central assumptions of that picture.

\section{Cycle solutions, not fixed points}

We tested 1000 Haar-random unitaries and 1000 Clifford-structured
unitaries on $(d_{\text{CR}}=2, d_{\text{CTC}}=3)$ for
\emph{cycle structure}: starting from a random pure initial state
$\rho_0$, we iterated $\rho_{n+1} = \mathcal{E}(\rho_n)$ and checked
whether the trajectory $\{\rho_n\}$ matched a previous iterate within
tolerance $10^{-8}$ within 100 steps.

\subsection*{Result}

\begin{table}[h]
\centering
\caption{Cycle-length distribution from 1000-channel sample.}
\label{tab:cycle-dist}
\begin{tabular}{lcc}
\toprule
Cycle length & Haar & Clifford-like \\
\midrule
1 (converged) & 65.9\% & \textbf{32.7\%} \\
2 & 17.5\% & \textbf{39.4\%} \\
3 & 5.8\% & 14.4\% \\
4 & 0.3\% & 10.3\% \\
$\ge 5$ or no cycle in 100 steps & 10.5\% & 0.0\% \\
\bottomrule
\end{tabular}
\end{table}

\textbf{The majority of Clifford-class D-CTC channels reach a limit
cycle, not a fixed point.} For Haar channels, fixed-point convergence
dominates (66\%) but cycle solutions are still present (24\%
combined).

\subsection*{Spectral prediction}

Channels with cycle length $k$ have second-largest eigenvalue
$\lambda_2(\mathbf{E})$ on the unit circle at angle $2\pi/k$:

\begin{center}
\begin{tabular}{lcc}
\toprule
Cycle length & Mean $|\arg \lambda_2|$ & Theoretical $2\pi/k$ \\
\midrule
1 (converged) & 0.000 & 0 \\
2 & 2.447 & $\pi \approx 3.142$ \\
3 & 2.094 & $2\pi/3 \approx 2.094$ \\
\bottomrule
\end{tabular}
\end{center}

The period-3 prediction is essentially exact. The period-2 cases
cluster near (but not at) $\pi$, suggesting a slight tilt off the
negative real axis.

Critically, \textbf{period-2 channels have $|\lambda_2| = 1.000$
exactly}, indicating a marginal mode on the unit circle. The
iteration's subdominant mode is neither decaying (which would give
convergence) nor growing (which would give instability) --- it
rotates in place. This marginal mode \emph{preserves input
information indefinitely}, which is the structural mechanism for
state-distinguisher amplification.

\subsection*{Cycles amplify more than fixed points}

Comparing amplification of two close mixed states (input trace
distance $0.1$) across cycle classes:

\begin{itemize}
\item Period-1 (convergent) channels: mean amplification 0.21
\item Period-2 (cycling) channels: mean amplification 0.43
\end{itemize}

Cycle solutions are roughly $2\times$ better state distinguishers
than fixed-point solutions, consistent with their marginal-mode
information preservation.

\section{Multi-basin dynamics: ergodicity overturned}

In a preliminary investigation we reported ``100\% ergodicity'' for
D-CTC iteration: every channel tested had a unique fixed point
regardless of initial state. That finding was an artifact of
restricting the initial-state ensemble to random pure states, which
mostly fall in the same basin of attraction.

\subsection*{Corrected test}

We probed each channel with 30 \emph{diverse} initial conditions:
random pure states, random mixed states (Dirichlet-weighted
superpositions of random pure states), basis states, equal-weight
superpositions, and the maximally-mixed state. We then clustered the
final iterates within Frobenius tolerance $10^{-7}$.

\begin{table}[h]
\centering
\caption{Number of distinct limit states per channel from 30 diverse
inits, $n = 200$ channels per ensemble.}
\label{tab:multi-basin}
\begin{tabular}{lcc}
\toprule
Ensemble & Mean distinct limits & Strict-unique fraction \\
\midrule
Haar & 1.0 & 100\% \\
Clifford-like & \textbf{24.78} & \textbf{18\%} \\
\bottomrule
\end{tabular}
\end{table}

A typical Clifford-class channel has approximately 25 distinct
limits, depending on initial state. Within a single cycle, however,
the iteration is ergodic: from 20 random inits all reaching a
particular period-2 channel, both phases of the cycle were
visited 10 times each.

\subsection*{Adversarial probing of one channel}

For a single Clifford-class channel we tested seven canonical
initial states (basis states, pairwise-equal superpositions, and
the maximally-mixed state) and observed at least four distinct
limit states with purities $\{1.0, 1.0, 0.5, 0.56, \ldots\}$.
\textbf{The same channel produces qualitatively different output
states for different input states} --- not just different
amplifications of the same fixed-point map, but genuinely distinct
limits.

\section{The encoding continuum}

The companion paper \cite{Knopp2026DCTC} distinguishes three
encoding regimes:
\begin{itemize}
\item \textbf{Abstract}: Clifford/Floquet/generic Hamiltonian $\to$
amplification independent of physical CTC content.
\item \textbf{Direct-CTC}: matrix entries literally proportional to
the spacetime CTC tensor $L_{\text{pair}}(r)$ $\to$ amplification
vanishes at chronology-protected configurations.
\item \textbf{LP-Hamiltonian}: physics-grounded Hamiltonian on a
Clifford backbone $\to$ pure but quiet (purity 1.0, amplification 0.18).
\end{itemize}

We test whether this trichotomy is a sharp phase boundary or a
continuum by constructing hybrid unitaries
\begin{equation}
U(\delta, t) = \mathrm{polar}\!\big[ (1-t) U_{\text{Clifford}} + t \cdot U_{\text{direct-L}}(\delta) \big]
\end{equation}
where polar denotes the polar decomposition projecting onto the
unitary group. We then sweep $\delta \in [0, 2\pi]$ and measure the
ratio of D-CTC amplification at $\delta = \pi$ (chronology-protected)
to amplification at $\delta = 0$ (maximum-CTC).

\begin{table}[h]
\centering
\caption{Smooth transition from abstract to direct-CTC encoding.}
\label{tab:smooth-transition}
\begin{tabular}{ccccc}
\toprule
$t$ & Amp at $\delta=0$ & Amp at $\delta=\pi$ & Ratio & Regime \\
\midrule
0.0 & 0.268 & 0.268 & 1.00 & abstract \\
0.5 & 0.310 & 0.140 & 0.45 & threshold \\
1.0 & 0.598 & 0.000 & 0.00 & direct-CTC \\
\bottomrule
\end{tabular}
\end{table}

The chronology-protection signature emerges gradually: at $t = 0$
the ratio is 1 (no signature), at $t = 0.5$ the ratio drops below
0.5 (signature threshold), and at $t = 1$ the ratio is 0 (full
signature). \textbf{No sharp phase boundary appears.} The trichotomy
of encoding regimes is a continuum parametrised by the channel's
direct-CTC content, and the 50\% mixing fraction marks the onset of
observable chronology-protection coupling.

\section{Re-interpretation of the D-CTC equation}

These three findings together suggest a substantial revision of the
standard reading of (\ref{eq:dctc-fp}). We propose that Deutsch's
self-consistency equation should be re-read as admitting
\emph{cycle solutions} of any period $k \ge 1$:
\begin{equation}
\rho_n = \rho_{n + k}, \quad k = 1, 2, 3, 4, \ldots
\end{equation}
with $k = 1$ being the special case of strict fixed points. Cycle
solutions of period $k \ge 2$:
\begin{itemize}
\item Exist generically in Clifford-class D-CTC channels (66\%
incidence at $d=6$);
\item Have $\lambda_2(\mathbf{E})$ exactly on the unit circle at
angle $2\pi/k$;
\item Are quantum-information-preserving (marginal mode);
\item Amplify state distinguishability roughly $2\times$ more than
period-1 fixed points;
\item Live in multi-basin dynamics where initial state selects the
particular cycle reached.
\end{itemize}

The Aaronson--Watrous polynomial-time PSPACE speedup
\cite{AaronsonWatrous2009} should therefore be re-interpreted as
operating on the cycle-solution class, not the fixed-point class.
This is consistent with the original theorem (which only requires
self-consistency, not strict convergence), and is sharper:

\begin{itemize}
\item \emph{The amplification mechanism is marginal-mode preservation},
not fixed-point projection.
\item \emph{Implementation requires basin-selection}: the protocol
must specify which limit cycle is reached, since multiple are
generically available.
\item \emph{Numerical practice} should detect cycles, not just
fixed points: a tolerance-based fixed-point solver will report
``convergence'' for a period-2 cycle because consecutive iterates
are within tolerance of each other, even though they alternate.
\end{itemize}

\section{Implications}

\subsection*{For computational complexity}

The Aaronson--Watrous result remains intact, but the channel class
realising it shifts: cycle-supporting channels (especially period-2)
are the natural site of polynomial-time PSPACE-style amplification.
Implementations via post-selection quantum circuits should
deliberately target cycle solutions rather than fixed points.

\subsection*{For physics interpretation}

If Deutsch's equation is read literally as a constraint on the
quantum state ``at the closed timelike curve'', cycle solutions
say the state need not be unique --- it can be a periodic
trajectory through finitely many density matrices. This is a
sharper statement than ``CTCs admit fixed points''.

\subsection*{For numerical practice}

Standard D-CTC simulators that use a fixed-point convergence
criterion will report ``convergence'' for cycle solutions of any
period as long as the convergence tolerance exceeds the
state-to-state distance within the cycle. We recommend always
checking explicitly for cycles by retaining the last $k = 1, 2, ...$
iterates and testing pairwise distances. The Systrophe library
includes such a cycle detector as \texttt{detect\_cycle\_length} in
the deep-dive scripts.

\section{Reproduction}

All findings are reproducible from:

\begin{verbatim}
python examples/dctc_deep_E1_cycles.py        # cycle distribution
python examples/dctc_deep_E4_trichotomy.py    # smooth transition
python examples/dctc_deep_E8_ergodicity.py    # multi-basin probe
\end{verbatim}

at \texttt{github.com/Zynerji/systrophe} v0.17.0 or later.

\section*{Acknowledgements}

These findings emerged from sustained empirical exploration with
the Anthropic Claude coding assistant. The author conceived the
three deep-dive experiments, the assistant implemented and executed
them, and we jointly interpreted the results. All numerical claims
are reproducible from the published scripts.

\bibliographystyle{plain}
\begin{thebibliography}{9}

\bibitem{Deutsch1991}
D.~Deutsch.
\newblock Quantum mechanics near closed timelike curves.
\newblock {\em Phys.\ Rev.\ D}, 44:3197--3217, 1991.

\bibitem{AaronsonWatrous2009}
S.~Aaronson and J.~Watrous.
\newblock Closed timelike curves make quantum and classical computing
equivalent.
\newblock {\em Proc.\ Roy.\ Soc.\ A}, 465:631--647, 2009.

\bibitem{Hawking1992}
S.~W.~Hawking.
\newblock Chronology protection conjecture.
\newblock {\em Phys.\ Rev.\ D}, 46:603--611, 1992.

\bibitem{Helstrom1976}
C.~W.~Helstrom.
\newblock {\em Quantum Detection and Estimation Theory}.
\newblock Academic Press, 1976.

\bibitem{Knopp2026DCTC}
C.~Knopp.
\newblock Clifford-structured Deutsch-CTC channels: empirical
amplification, spectral oracle, and the encoding-dependence of
chronology protection.
\newblock \url{https://github.com/Zynerji/systrophe/blob/main/paper/dctc_amplification.pdf}, 2026.

\bibitem{Knopp2026Systrophe}
C.~Knopp.
\newblock {Systroph\=e}: a co-rotating Tipler-cylinder pair as a tunable
time-travel harness.
\newblock \url{https://github.com/Zynerji/systrophe}, 2026.

\end{thebibliography}

\end{document}
